% Generated draft: requires mathematical and source review.
\documentclass[11pt]{article}
\usepackage[margin=25mm]{geometry}
\usepackage{fontspec}
\setmainfont{Latin Modern Roman}
\usepackage{amsmath,amssymb,mathtools,unicode-math}
\setmathfont{Latin Modern Math}
\usepackage{enumitem,needspace,xurl,hyperref}
\setlength{\parindent}{0pt}
\setlength{\parskip}{5pt}
\setlength{\emergencystretch}{3em}
\newcommand{\N}{\mathbb{N}}
\newcommand{\Z}{\mathbb{Z}}
\newcommand{\Q}{\mathbb{Q}}
\newcommand{\R}{\mathbb{R}}
\newcommand{\C}{\mathbb{C}}
\newcommand{\F}{\mathbb{F}}
\newcommand{\A}{\mathbb{A}}
\newcommand{\PP}{\mathbb{P}}
\newcommand{\E}{\mathbb{E}}
\newcommand{\eps}{\varepsilon}
\newcommand{\dd}{\,\mathrm{d}}
\begin{document}
\section*{UnsolvedMath ID: 30004405}
\textbf{Source problem code:} catalogue entry
\subsection{Definitions and mathematical statement}
Write \(\Gamma_K=\pi_1(S^3\setminus K,x_0)\), with a fixed basepoint \(x_0\) in the complement. The supplied statement does not specify the admissible objects \(K\); its title suggests hyperbolic knots, but this hypothesis is absent from the displayed statement. All notation below is conditional on specifying \(K\).

For a subgroup \(H\leq\Gamma_K\) and an element \(\sigma\in\Gamma_K\setminus H\), define the separating index by
\[
\operatorname{Sep}_{\Gamma_K}(H,\sigma)
=
\min\bigl\{[\Gamma_K:G]: H\leq G\leq\Gamma_K,\ \sigma\notin G,\ [\Gamma_K:G]<\infty\bigr\},
\]
with value \(+\infty\) when the displayed set is empty. Here \([\Gamma_K:G]\) is the number of left cosets of \(G\) in \(\Gamma_K\). The subgroup inclusions are not required to be proper.

The source asks: given \(H\leq\Gamma_K\) and \(\sigma\in\Gamma_K\setminus H\), can one give a finite quantitative upper bound for \(\operatorname{Sep}_{\Gamma_K}(H,\sigma)\)? Equivalently, can one bound the index of a finite-index subgroup \(G\leq\Gamma_K\) satisfying
\[
H\leq G,\qquad \sigma\notin G?
\]
No existence of such a finite-index subgroup is asserted here.

The displayed source formulation permits arbitrary subgroups \(H\), whereas its supplied literature assessment describes the remaining problem as concerning finitely generated subgroups. These give different quantifier scopes. A possible finitely generated reading restricts the question to \(H=\langle h_1,\ldots,h_r\rangle\) for some integer \(r\geq0\) and elements \(h_i\in\Gamma_K\), but that restriction is not present in the displayed statement and is not adopted silently. The record also does not specify the parameters on which the bound should depend, any generating set or word-length convention, a representation of the input subgroup, or a required growth rate. Consequently, the precise quantitative target and the intended hypotheses remain unresolved.
\subsection{Short English statement}
Given a subgroup of the intended knot or link complement group and an element outside it, seek a finite-index subgroup containing the first subgroup and excluding that element, together with a quantitative bound on its index. The admissible complements, whether the subgroup must be finitely generated, and the permitted dependence of the bound require clarification.
\subsection{Sources}
\begin{itemize}
\item Supplied local UnsolvedMath registry record, permanent ID \(30004405\), source problem code \texttt{catalogue entry}, titled \emph{Effective Subgroup Separability in Hyperbolic Knot Groups}; collection: Oberwolfach Reports Open Problems. The record gives proposed year 2020 and supplies no named proposer.
\item Supplied report citation: \emph{Low-dimensional Topology (2021)}, Oberwolfach Reports, DOI 10.4171/owr/2020/8, pp. 465--516. \url{https://doi.org/10.4171/owr/2020/8}. The literature list also labels this reference \emph{Oberwolfach Report 8/2020}. These supplied year labels are retained without bibliographic reconciliation.
\item Bering and Futer, \emph{Explicit special covers of alternating links}. \url{https://arxiv.org/abs/1909.12266}. This attribution and link are supplied by the record; the paper has not been independently checked here.
\item Supplied literature assessment, dated exactly \texttt{2026-08-22}: the record labels the problem partially solved and states that quantitative separability is established for closed curves and Abelian subgroups in prime alternating link complements, while the stated finitely generated subgroup problem remains unresolved. This is supplied dated curation derived from the OWR material, not an independently verified literature assessment. Its claims and time-sensitive status require independent verification.
\end{itemize}
\end{document}
